%% ============================================================
%% Chapter 1 — Introduction
%% ============================================================

\chapter{Introduction}\label{chapter:1_Introduction}

\section{Context}
\label{sec:intro:context}

Mandarin Chinese has approximately 1.18~billion native and second-language speakers worldwide, including about 194~million second-language speakers according to Ethnologue's global speaker estimates~\cite{ethnologue2025}. These figures help explain why Mandarin remains an important target language in education and professional communication. Unlike the majority of Indo-European languages, Mandarin is a \emph{tonal} language: the pitch contour realised at the syllable level is phonemically distinctive, meaning that a change in tone can change the meaning of a word entirely~\cite{jongman2006tones}. Standard Mandarin employs four lexical tones: a high-level tone~(T1), a rising tone~(T2), a dipping tone~(T3), and a falling tone~(T4), plus a reduced neutral tone, represented in this thesis as the fifth classification category, T5~\cite{zahner2022toneexperience}. For example, the syllable \textit{ma} can mean ``mother'' (T1), ``hemp'' (T2), ``horse'' (T3), or ``scold'' (T4), depending solely on its lexical tone~\cite{jongman2006tones}.

The tonal nature of Mandarin poses a substantial challenge for second-language~(L2) learners whose first language is non-tonal. In non-tonal languages such as Polish, English, or French, pitch variation serves only a prosodic function, signalling questions, emphasis, or discourse boundaries, rather than a lexical one. Consequently, learners from these linguistic backgrounds tend to underperceive tonal contrasts and produce pitch contours that deviate systematically from native targets~\cite{hao2012}. Research on L2 tone acquisition has identified certain tone pairs as particularly problematic: the T2/T3 contrast, where both tones begin in the low-pitch region, is consistently the most confusable for non-tonal-L1 learners~\cite{hao2012}. A recent study specifically targeting Polish learners of Mandarin reported systematic differences in both pitch height and contour shape relative to native speakers across disyllabic tonal combinations~\cite{chu2024statistical}. Taken together, these studies indicate that Polish learners constitute a relevant population for evaluating automatic methods intended to support Mandarin tone training.

Computer-Assisted Pronunciation Training~(CAPT) systems aim to address this need by providing automated, scalable feedback on learner pronunciation~\cite{neri2002,rogerson2021}. A CAPT system typically comprises an automatic speech recognition or classification component that analyses the learner's production, coupled with a feedback mechanism that communicates results. In the specific domain of Mandarin tone assessment, early classifiers adopted convolutional neural network~(CNN) architectures that operate on pitch-contour features~\cite{chen2016cnn} or mel-spectrogram representations~\cite{gao2019tonenet}, achieving high accuracy on isolated syllables under controlled conditions. While effective, these models rely on hand-crafted acoustic features and are trained from scratch, limiting their ability to generalise across speakers, recording conditions, and levels of proficiency.

Recent work in speech processing has increasingly relied on self-supervised learning~(SSL) models such as Wav2Vec~2.0~\cite{baevski2020} and HuBERT~\cite{hsu2021}. These models are pre-trained on large amounts of unlabelled speech and can then be adapted to downstream tasks with comparatively limited labelled data. In Mandarin speech research, prior studies have shown that intermediate SSL representations retain information relevant to lexical tone and other suprasegmental categories, and that task-specific fine-tuning can improve tone-related analysis~\cite{yuan2023,delafuente2024}. For this reason, SSL-based classifiers are a plausible alternative to earlier CNN-based tone-classification systems, especially in settings where learner speech is variable and labelled data are limited.

\section{Identification of Gaps}
\label{sec:intro:gaps}

The literature reviewed above leaves three issues unresolved for the present study.

\textbf{Population gap.} Most tone-classification systems discussed in this thesis were developed and evaluated on native Mandarin speech~\cite{chen2016cnn,gao2019tonenet}. Publicly described learner-oriented resources remain limited, and the non-native corpus most directly relevant to pronunciation assessment in the present literature is OMPAL, which contains speech from French~L1 learners rather than Polish learners~\cite{hsieh2025ompal}. For Polish learners of Mandarin, the available studies are acoustic and pedagogical in focus: Chu et~al.~\cite{chu2024statistical} provide a quantitative analysis of disyllabic tone production, and Zajdler~\cite{zajdler2021} examines tonal production in read speech, but neither study develops or evaluates an automatic tone-classification model. To the best of the author's knowledge, no published study has evaluated modern Mandarin tone classifiers on speech produced by Polish learners.

\textbf{Model comparison gap.} While both CNN-based and SSL-based approaches have been applied to Mandarin tone classification independently, no systematic comparison of these paradigms has been conducted on non-native learner speech. Furthermore, SSL models can be pre-trained on different data, including native Mandarin corpora, mixed native-and-learner corpora, or multilingual checkpoints, and the effect of these pre-training strategies on downstream tone classification accuracy for L2 speakers remains unexplored.

\textbf{Expert alignment gap.} Existing tone classification models are typically evaluated against transcription labels that indicate the intended tone. However, in a CAPT context, what matters is whether the model's assessment agrees with how a human expert perceives the learner's pronunciation quality. There is a lack of evaluation frameworks that align model predictions with expert phonetic assessments, making it difficult to judge whether a classifier is genuinely useful as a pronunciation feedback tool.

\section{Goals of the Thesis}
\label{sec:intro:goals}

Accordingly, this thesis compares CNN-based and SSL-based models for Mandarin tone classification on recordings produced by Polish learners and examines how closely model predictions align with expert judgements. This objective is pursued through the following specific tasks:

\begin{enumerate}
  \item Compare six tone classification setups on recordings of Polish learners of Mandarin: a CNN baseline trained on mel-spectrograms, three Wav2Vec~2.0 variants differing in their initialisation and adaptation data (learner-only fine-tuning, AISHELL-1 native-speech adaptation~\cite{bu2017aishell}, and a Chinese-specific checkpoint), and two Chinese HuBERT variants differing in pooling and loss design.
  \item Determine which pre-training strategy, backbone, and training setup yield the highest classification accuracy on non-native tone productions.
  \item Evaluate the agreement between model predictions and expert phonetic assessments to establish whether the models' judgements align with qualified human evaluation.
  \item Provide practical guidance for integrating the best-performing model into a CAPT pipeline for Mandarin tone training.
\end{enumerate}

\section{Outline of the Thesis}
\label{sec:intro:outline}

The thesis is organised as follows:

\begin{itemize}
  \item Chapter~\ref{chapter:2_StateOfTheArt} reviews the relevant literature on Mandarin tone phonology, second-language tone acquisition, computer-assisted pronunciation training, and self-supervised speech modelling.
  \item Chapter~\ref{chapter:3_Proposal} defines the research problem, states the main assumptions and requirements, and specifies the acceptance criteria for the study.
  \item Chapter~\ref{chapter:4_Design} describes the system architecture, model configurations, and data-processing pipeline used in the implementation.
  \item Chapter~\ref{chapter:5_Experiments} presents the experimental setups, evaluation procedure, quantitative results, and comparative analysis.
  \item Chapter~\ref{chapter:6_Summary} summarises the main findings, discusses limitations, and outlines directions for future work.
\end{itemize}
